\documentclass[10pt,a4paper]{amsart}
\usepackage[margin=19mm]{geometry}
\usepackage{amsmath,amssymb,mathtools}
\usepackage[scr=rsfs,cal=euler]{mathalfa}
\usepackage{xfrac}
\usepackage[unicode,hidelinks,bookmarks=false]{hyperref}
\newcommand{\ie}{i.e.\ }
\newcommand{\cf}{cf.\ }
\newcommand{\resp}{respectively\ }
\newcommand{\Subsection}{\subsection}
\providecommand{\nobreakdash}{\nobreak\hbox{-}}
\setlength{\emergencystretch}{2em}
\multlinegap0pt
\usepackage{bm}
\usepackage{bbm}
\usepackage{dsfont}
\usepackage{xspace}
\usepackage{cancel}
\usepackage{mathtools}
\usepackage{amsthm}
\makeatletter
\@ifundefined{theorem}{\newtheorem{theorem}{Theorem}}{}
\makeatother
\title[The local antimagic (total) chromatic numbers of firecracker]{The local antimagic (total) chromatic numbers of firecracker graphs and edge-corona product graphs}
\author{Xue Yang, Hong Bian, Xueliang Li, Zhixia Yang, Haizheng Yu}
\date{Source: arXiv:2602.01562v1 (CC BY 4.0)}
\begin{document}
\maketitle

\noindent\emph{Source and licence.} The local antimagic (total) chromatic numbers of firecracker graphs and edge-corona product graphs. Version arXiv:2602.01562v1 (CC BY 4.0), distributed under the Creative Commons Attribution 4.0 International licence, as recorded in the arXiv licence metadata for this exact version and shown on the abstract page https://arxiv.org/abs/2602.01562v1. Full rights notice: licenses/arxiv-2602.01562-CC-BY-4.0.txt.\par\smallskip

\noindent\textit{Editorial scope.} Counted unit: the Theorem whose source label is \texttt{None} in the article (source0.tex lines 605--610, source label \texttt{None}), delivered together with the article's own complete original proof of it (source0.tex lines 612--827, one \texttt{proof} environment of 216 lines). No mathematics is written, repaired or re-derived: statement and proof are the article's own text line for line. Required context retained uncounted: Theorem 2.3 (lines 469--474). Every definition, display and label of those lines is retained unchanged. Omitted from this package and not redistributed: 1--468, 475--604, 611--611, 828--1736. Nothing inside a retained excerpt is omitted or altered: every retained body is delivered exactly as the article's lines are written, so no omission receipt of any kind accompanies this package. Citations: the retained excerpts carry no citation command at all, so no bibliography entry of the article is transcribed and no reference is redistributed. Reference adaptation: none was needed and none was made; every reference inside the delivered unit resolves inside it, so no unresolved reference or citation is produced. Printed slips kept literal: no printed word, symbol, hypothesis or number of the counted statement or of its proof was altered. Layout and class: the article's own class is \texttt{article} with packages including geometry, graphicx, diagbox, color, cite, array, amsthm, amsmath, amssymb, indentfirst, arydshln, rotating, dashrule, epstopdf; the wrapper is the lane's pinned \texttt{amsart} editorial wrapper at 10pt on A4, which restates the theorem-like environments the retained text needs and adds the article's own macro definitions that the retained text needs (none), with extra wrapper packages (bm, bbm, dsfont, xspace, cancel, mathtools). The delivered numbering is the unit's own. Assets: no retained span contains a figure environment, a graphics-inclusion command, a PDF-page inclusion, a table or a TikZ picture; no image file is copied into this package, the package declares no asset dependency and no image file is redistributed with it. Licence: version arXiv:2602.01562v1 of this article is distributed under the Creative Commons Attribution 4.0 International licence, as recorded in the arXiv licence metadata for this exact version and shown on the abstract page https://arxiv.org/abs/2602.01562v1. A full rights notice is delivered at licenses/arxiv-2602.01562-CC-BY-4.0.txt. Characters: every retained excerpt is pure ASCII, so the delivered engine, XeLaTeX with its default Latin Modern text font, reports no missing character.
\begin{theorem}\label{Theorem 2.3}
For firecracker graph $F_{n,k}$, the lower bound of the local antimagic total chromatic number satisfies the following inequality
\begin{equation*}
  \chi_{lat}(F_{n,k})\geq \chi(F_{n,k})=2.
\end{equation*}
\end{theorem}

\begin{theorem}
  For $n\geq 3$ and $k\geq 2$, the local antimagic total chromatic number of the firecracker graph $F_{n,k}$ is bounded as follows:
\begin{equation*}
  2\leq \chi_{lat}(F_{n,k})\leq 3.
\end{equation*}
\end{theorem}

\begin{proof}
From Theorem \ref{Theorem 2.3}, the lower bound of the local antimagic total chromatic number of $F_{n,k}$ for $n\geq 3$ and $k\geq 2$ is given by $\chi_{lat}(F_{n,k})\geq\chi(F_{n,k})=2$.

The upper bound of the local antimagic total chromatic number of $F_{n,k}$ is determined by the colors induced by the local antimagic total labeling. Define a bijection $g:V(F_{n,k})\cup E(F_{n,k})\rightarrow [1,2nk+2n-1]$, and let $U=[1,2nk+2n-1]$.
Discuss the local antimagic total labeling of graph $F_{n,k}$ through the following cases.

Case 1. $k$ is even.

Initially, we utilize the set $A=[2n,nk-1]\cup[nk+3n,2nk+n-1]$ to label the edges and vertices of $F_{n,k}$ excluding its subgraph $F_{n,2}$.

$g(u_{i}v_{i,j})=\begin{cases}
nj-n-1+i, & \text{if $1\leq i\leq n$, $3\leq j\leq k$ and odd $j$},\\[4pt]
nj-i,        & \text{if $1\leq i\leq n$, $3\leq j\leq k$ and even $j$}.\end{cases}$\vspace{0.4em}

$g(v_{i,j})=\begin{cases}
2nk+4n-nj-i, & \text{if $1\leq i\leq n$, $3\leq j\leq k$ and odd $j$},\\[4pt]
2nk+3n-1-nj+i,        & \text{if $1\leq i\leq n$, $3\leq j\leq k$ and even $j$}.\end{cases}$\vspace{0.6em}

Subsequently, employ the set $B=U\setminus A$ to label the subgraph $F_{n,2}$ of $F_{n,k}$, where $B=[1,2n-1]\cup[nk,nk+3n-1]\cup[2nk+n,2nk+2n-1]$.
It is further divided into the following two subcases according to the parity of $n$.

Subcase 1.1.  $n$ is odd.

When $n=3$, the local antimagic total labeling of $F_{3,k}$ for even $k$ is specified as follows:

$g(v_{1,1}v_{2,1})=1,\hspace{0.5em} g(v_{2,1}v_{3,1})=2, \hspace{0.5em}g(u_{1}v_{1,2})=3, \hspace{0.5em}g(u_{2}v_{2,2})=4, \hspace{0.5em}g(u_{3}v_{3,2})=5,$

$g(u_{1}v_{1,1})=3k+5,\hspace{0.5em}g(u_{2}v_{2,1})=3k+3,\hspace{0.5em}g(u_{3}v_{3,1})=3k+6,$\vspace{0.5em}

$\hspace{-0.5em}\begin{array}{lll}
g(u_{1})=3k+7,&g(u_{2})=3k+8,&g(u_{3})=3k+4,\\[4pt]
g(v_{1,1})=3k+2,&g(v_{2,1})=3k+1,&g(u_{3}v_{3,1})=3k,\\[4pt]
g(v_{1,2})=6k+5,&g(v_{2,2})=6k+4,&g(u_{3}v_{3,2})=6k+3,
\end{array}$\vspace{0.6em}

For $n\geq 5$, the corresponding local antimagic total labeling of $F_{n,k}$ for odd $n$ and even $k$ is obtained by\\

$g(v_{i,1}v_{i+1,1})=\begin{cases}
i, & \text{if $1\leq i<n$ and odd $i$},\\
n+1-i, &\text{if $1\leq i<n$ and even $i$}.\end{cases}$\vspace{0.5em}

$g(u_{i}v_{i,1})=\begin{cases}
nk+2n-1, & \text{if $ i=1$},\\
\frac{2nk+4n-3-i}{2},& \text{if $2\leq i< n$ and odd $i$},\\[4pt]
\frac{2nk+3n-1-i}{2}, &\text{if $2\leq i\leq n$ and even $i$},\\[4pt]
\frac{2nk+5n-3}{2}, & \text{if $i=n$}.\end{cases}$\vspace{0.5em}

$g(u_{i}v_{i,2})=\begin{cases}
n, & \text{if $ i=1$},\\
n+1+i, & \text{if $2\leq i< n$ and odd $i$},\\
n-1+i,  &\text{if $2\leq i\leq n$ and even $i$},\\
n+2, & \text{if $ i=n$}.\end{cases}$\vspace{0.5em}

$g(u_{i})=\begin{cases}
\frac{2nk+5n-1}{2}, & \text{if $ i=1$},\\[4pt]
\frac{2nk+5n-2-i}{2},& \text{if $2\leq i< n$ and odd $i$},\\[4pt]
\frac{2nk+6n-i}{2}, &\text{if $2\leq i\leq n$ and even $i$},\\[4pt]
nk+2n-2, & \text{if $ i=n$}.\end{cases}$\vspace{0.5em}

$g(v_{i,1})=\begin{cases}
nk+n-1, & \text{if $ i=1$},\\
\frac{2nk-3+i}{2}, & \text{if $2\leq i\leq n$ and odd $i$},\\[4pt]
\frac{2nk+n-3+i}{2}, &\text{if $2\leq i\leq n$ and even $i$}.\end{cases}$\vspace{0.5em}

$g(v_{i,2})=\begin{cases}
2nk+2n-1, & \text{if $ i=1$},\\
2nk+2n-2-i, & \text{if $2\leq i< n$ and odd $i$},\\[4pt]
2nk+2n-i,  &\text{if $2\leq i\leq n$ and even $i$},\\[4pt]
2nk+2n-3, & \text{if $ i=n$}.\end{cases}$\vspace{0.6em}

In this subcase, it can be calculated from the above labels that

$\omega_{t}(v_{s,1})=\omega_{t}(v_{i,j})=2nk+3n-1$ for $1\leq s,i\leq n$, $2\leq j\leq k$ and odd $s$,

$\omega(v_{s,1})=2nk+3n-2$ for $1\leq s\leq n$ and even $s$,

$\omega_{t}(u_{i})=\frac{nk^{2}+4nk+7n-k-1}{2}$ for $1\leq i\leq n$.

Clearly, this local antimagic total labeling of $F_{n,k}$ induces three different colors, and so $\chi_{lat}(F_{n,k})\leq 3$ for odd $n$ and even $k$.

Subcase 1.2.  $n$ is even.

The local antimagic total labeling for even $n$ is defined as follows:\\

\vspace{-0.8em}
$g(v_{i,1}v_{i+1,1})=\begin{cases}
                       \frac{i+1}{2}, & \text{if $1\leq i<n$ and odd $i$},\\[4pt]
                       \frac{n+i}{2}, & \text{if $1\leq i<n$ and even $i$}.\end{cases}$\vspace{0.5em}

$g(u_{i}v_{i,1})=\begin{cases}
nk+n-\frac{i+1}{2}, & \text{if $1\leq i\leq n-2$ and odd $i$},\\[4pt]
nk-1+\frac{n-i}{2}, & \text{if $1\leq i\leq n-2$ and even $i$},\\[4pt]
nk+n,     &\text{if }i=n-1,\\[4pt]
nk-1+\frac{n}{2}, &\text{if }i=n.\end{cases}$\vspace{0.4em}

$g(u_{i}v_{i,2})=\begin{cases}
\frac{3n-3-i}{2}, & \text{if $1\leq i\leq n-2$ and odd $i$},\\[4pt]
\frac{4n-2-i}{2}, & \text{if $1\leq i\leq n-2$ and even $i$},\\[4pt]
\frac{3n-2}{2},     &\text{if } i=n-1,\\[4pt]
2n-1, &\text{if }i=n.\end{cases}$\vspace{0.4em}

$g(u_{i})=\begin{cases}
nk+2n+1+i, & \text{if $1\leq i\leq n-2$},\\[4pt]
nk+n+1+i,        & \text{if $i=n-1,n$}.\end{cases}$\vspace{0.4em}

$g(v_{i,1})=\begin{cases}
nk+2n-1, & i=1,\\[4pt]
nk-1+\frac{3n-i}{2}, & \text{if $1<i\leq n-2$ and odd $i$},\\[4pt]
nk+2n-2-\frac{i}{2},        & \text{if $1< i\leq n-2$ and even $i$},\\[4pt]
nk+\frac{n}{2},  &\text{if } i=n-1,\\[4pt]
nk+2n-2, &\text{if }i=n.\end{cases}$\vspace{0.4em}

$g(v_{i,2})=\begin{cases}
2nk+\frac{3n+1+i}{2}, & \text{if $1\leq i\leq n-2$ and odd $i$},\\[4pt]
2nk+n+\frac{i}{2},        & \text{if $1\leq i\leq n-2$ and even $i$},\\[4pt]
2nk+\frac{3n}{2},         & \text{if } i=n-1,\\[4pt]
2nk+n, &\text{if }i=n.\end{cases}$\vspace{0.4em}

\noindent According to the above labels, we have

$\omega_{t}(v_{s,1})=\omega_{t}(v_{i,j})=2nk+3n-1$ for $1\leq s,i\leq n$, $2\leq j\leq k$ and odd $s$,

$\omega_{t}(v_{s,1})=2nk+3n-3$ for $1\leq s\leq n$ and even $s$,

$\omega_{t}(u_{i})=\frac{nk^{2}+4nk+5n-k}{2}+2nk+2n$ for $1\leq i\leq n$.

\noindent In this subcase, the local antimagic total labeling of $F_{n,k}$ uses three distinct colors, and thus $\chi_{lat}(F_{n,k})\leq 3$ for even $n$ and $k$.

Therefore, it is verified that $\chi_{lat}(F_{n,k})\leq 3$ for even $k$.

Case 2. $k$ is odd.

Similarly, use the set $C=[n,nk-1]\cup[nk+3n,2nk+2n-1]$ to label the edges and vertices of the firecracker graph $F_{n,k}$ excluding its subgraph $F_{n,1}$, in the following way:\\
\vspace{0.5em}

$g(u_{i}v_{i,j})=\begin{cases}
     nj-n-1+i, & \text{for $1\leq i\leq n$, $2\leq j\leq k$, and even $j$}, \\
     nj-i, & \text{for $1\leq i\leq n$, $2\leq j\leq k$, and odd $j$}. \end{cases}$\vspace{0.5em}

$ g(v_{i,j})=\begin{cases}
     2nk+4n-nj-i, & \text{for $1\leq i\leq n$, $2\leq j\leq k$, and even $j$}, \\
     2nk+3n-1-nj+i, & \text{for $1\leq i\leq n$, $2\leq j\leq k$, and odd $j$}. \end{cases}$\vspace{0.5em}

Then, use the set $D=U\setminus C=[nk,nk+3n-1]$ to label the subgraph $F_{n,1}$ of $F_{n,k}$. We discuss the parity of $n$ in a manner consistent with previous discussion.

Subcase 2.1. $n$ is odd.

For $n=3$,  present the local antimagic toal labeling of $F_{3,k}$ by the following formulas:

$g(v_{1,1}v_{2,1})=1,\hspace{1em}g(v_{2,1}v_{3,1})=2,\hspace{1em}g(u_{1}v_{1,1})=3k+7,\hspace{1em}g(u_{2}v_{2,1})=3k+8$

$g(u_{3}v_{3,1})=3k+5,\hspace{1em}g(v_{1,1})=3k,\hspace{1em}g(v_{2,1})=3k+2,\hspace{1em}g(v_{3,1})=3k+1$,

$g(u_{1})=3k+4,\hspace{1em}g(u_{2})=3k+3,\hspace{1em}g(u_{3})=3k+6$.

When $n\geq 5$, give the local antimagic toal labeling of $F_{n,k}$ by\vspace{0.6em}

$g(v_{i,1}v_{i+1,1})=\begin{cases}
\frac{2n-1-i}{2}, & \text{if $1\leq i\leq n$ and odd $i$},\\[4pt]
\frac{i}{2}, & \text{if $1\leq i\leq n$ and even $i$}. \end{cases}$\vspace{0.5em}

$g(u_{i}v_{i,1})=\begin{cases}
\frac{2nk+4n-1-i}{2}, & \text{if $1\leq i\leq n-1$ and odd $i$},\\[4pt]
\frac{2nk+6n-i}{2}, & \text{if $1\leq i\leq n-1$ and even $i$},\\[4pt]
\frac{2nk+5n-1}{2}, & \text{if $ i= n$}.\end{cases}$\vspace{0.5em}

$g(v_{i,1})=\begin{cases}
\frac{2nk+1+i}{2}, & \text{if $1\leq i\leq n-1$ and odd $i$},\\[4pt]
\frac{2nk+n-1+i}{2}, & \text{if $1\leq i\leq n-1$ and even $i$},\\[4pt]
nk, & \text{if $ i= n$}.\end{cases}$\vspace{0.5em}

$g(u_{i})=\begin{cases}
\frac{2nk+4n-1+i}{2}, & \text{if $1\leq i\leq n-1$ and odd $i$},\\[4pt]
\frac{2nk+2n-2+i}{2}, & \text{if $1\leq i\leq n-1$ and even $i$},\\[4pt]
\frac{2nk+3n-1}{2}, & \text{if $ i= n$}.\end{cases}$\vspace{0.6em}

\noindent By the above labels, we obtain that

$\omega_{t}(v_{s,1})=\omega_{t}(v_{i,j})=2nk+3n-1$ for $1\leq s,i\leq n$, $2\leq j\leq k$ and odd $s$,

$\omega_{t}(v_{s,1})=\frac{4nk+9n-1}{2}$ for $1\leq s\leq n$ and even $s$,

$\omega_{t}(u_{i})=\frac{nk^{2}+4nk+7n-k-1}{2}$ for $1\leq i\leq n$.

\noindent These are three distinct colors, and so $\chi_{lat}(F_{n,k})\leq 3$ when $n$ and $k$ are odd.

Subcase 2.2.  $n$ is even. \vspace{0.4em}

We give the following labeling\\

$g(v_{i,1}v_{i+1,1})=\begin{cases}
                       \frac{2n-1-i}{2}, & \text{if $1\leq i<n$ and odd $i$},\\[4pt]
                       \frac{i}{2}, & \text{if $1\leq i<n$ and even $i$}.\end{cases}$\vspace{0.4em}

$ g(u_{i}v_{i,1})=\begin{cases}
           \frac{2nk+3n-1-i}{2}, & \text{if $1\leq i\leq n$ and odd $i$},\\[4pt]
           \frac{2nk+4n-i}{2}, & \text{if $1\leq i\leq n$ and even $i$}.\end{cases}$\vspace{0.4em}

$g(u_{i})=\begin{cases}
    \frac{2nk+5n-1+i}{2}, & \text{if $1\leq i\leq n$ and odd $i$},\\[4pt]
    \frac{2nk+4n-2+i}{2}, & \text{if $1\leq i\leq n$ and even $i$}.\end{cases}$\vspace{0.4em}

$g(v_{i,1})=\begin{cases}
           \frac{2nk+n-3+i}{2}, & \text{if $1\leq i<n$ and odd $i$},\\[4pt]
            \frac{2nk-2+i}{2},    & \text{if $1\leq i<n$ and even $i$},\\[4pt]
                       nk+n-1, &\text{if } i=n.\end{cases}$\vspace{0.4em}

\noindent From the above labels in this subcase, it can be calculated that

$\omega_{t}(v_{s,1})=\omega_{t}(v_{i,j})=2nk+3n-1$ for $1\leq s,i\leq n$, $2\leq j\leq k$ and even $s$,

$\omega_{t}(v_{s,1})=2nk+3n-3$ for $1\leq s\leq n$ and odd $s$,

$\omega_{t}(u_{i})=\frac{nk^{2}+4nk+7n-k-1}{2}$ for $1\leq i\leq n$.

\noindent Obviously, there are three different colors used by the local antimagic total labeling, and so $\chi_{lat}(F_{n,k})\leq 3$ for even $n$ and odd $k$. In this case, it is concluded that $\chi_{lat}(F_{n,k})\leq 3$ for odd $k$. The conclusion is thus proved.\end{proof}

\end{document}
